\DocumentMetadata{uncompress,lang=en,testphase={phase-III}}
\documentclass{article}
\usepackage[lazy,gb4e]{linguexx}
\begin{document}

\GlossTierLang{1}{de}
% the free translation gets its own language, distinct from the object
% tier (de) and from the document (en), so /Lang (fr) can only come
% from \GlossTransLang reaching the \glt paragraph
\GlossTransLang{fr}
% PDF tagging smoke test: every construct here has, at some point,
% corrupted the structure tree (mismatched text-unit bookkeeping) under
% the latex-lab testphase code.  -halt-on-error turns any tagpdf
% accounting error into a failed compile, which is the assertion that
% matters; the text checks below just confirm the content arrived.
\noindent CCC margin reference.

AAA start.

\ex. MAINTEXT one.
\a. LETTERA plain sub.
\b. *LETTERB judged sub.
\a. ROMANI roman sub.
\z.\z.

\exg. GLOSSA GLOSSB GLOSSC.\\
      glossa glossb \lpzg{3sg.pst}.\\
\glt `Translation.'

\ex. I saw \altn{the cat}{a dog} here.

\exg. ALTGHOST \altg{Socke}{Tonne} ALTGVERB.\\
      altghost \altg{sock.\lpzg{sg}}{ton.\lpzg{sg}} altgverb.\\

FNHOST here.\footnote{In the note: \ex. FNEX example.

}

\begin{exe}
\ex[??]{EXEITEM one.}
\end{exe}

BBB following paragraph after everything.

% The abbreviation list is a real tagged list; its labels must stay flush
% left (the block code re-boxes a short label flush right otherwise).
LPZGLIST follows.
\lpzglist[ignore={3}]
\end{document}
